\chapter{Case Study: Building a Lua-Preprocessed Assembler}
\label{ch:lua-assembler}

An assembler becomes a useful case study when its output survives another tool's interpretation.
Producing a plausible hex dump is insufficient: a symbol may name the wrong section, a relocation may point one byte too far, or a valid instruction may be assigned the wrong operand width.
The implementation in \path{example-code/lua-toolchain/luaas.c} therefore produces real x86-64 ELF64 objects and participates in two independently checked linking paths.
This chapter follows its data structures and decisions from a Lua source file to the final relocation records.

The implementation is C11 with the Lua 5.4 library.
It contains its own instruction encoder and ELF writer; it does not invoke another assembler to produce its output.
GNU \code{as}, GNU \code{ld}, and \code{readelf} are independent verification tools used by the test driver.
Their role must remain separate from the implementation being tested.

\section{Set an executable contract}
\label{sec:ch56-contract}

The command accepts a Lua source file and an output object path.
With \code{--asm}, it instead reads already expanded assembly text.
Both paths feed the same parser, encoder, and object writer.
This second entry path is valuable for isolating an encoding problem from a preprocessing problem: a failing generated line can be reproduced without executing its Lua generator.

The supported registers are the first eight x86 general-purpose registers, in 32-bit or 64-bit forms.
The instruction set includes immediate and register moves, register XOR, immediate add/subtract/compare, RIP-relative address formation, near call/jump/conditional branches, return, no-operation, undefined-instruction trap, and system call.
There is no general memory-operand grammar, SIMD family, extended-register encoding, or short-branch selection.
The purpose of a subset is to make every accepted form precise and tested.

Table~\ref{tab:ch56-contract} distinguishes implemented features from the design topics developed in the preceding chapters.
An unsupported mnemonic or directive must fail explicitly rather than being ignored as an annotation.
The writer does not publish a partial object when an earlier phase fails.

\begin{table}[htbp]
\caption{The case-study interface defines a coherent implemented subset rather than an approximate full assembler.}
\label{tab:ch56-contract}
\centering
\begin{tabular}{p{0.29\linewidth}p{0.58\linewidth}}
\toprule
Facility & Implemented contract \\
\midrule
Preprocessing & Lua 5.4 script with emit; optional plain-text input. \\
Sections & Text, read-only data, writable data, and NOBITS storage. \\
Symbols & Local labels, explicit global/external/weak declarations. \\
Expressions & Signed 64-bit integer or one symbol plus integer addend. \\
Relocations emitted & AMD64 64, PC32, and unsigned 32. \\
Layout & Fixed instruction forms and power-of-two alignment through 4096. \\
Not implemented & General expression algebra, short-branch relaxation, debug/unwind emission. \\
\bottomrule
\end{tabular}
\end{table}

\section{Keep ownership simple enough to inspect}
\label{sec:ch56-storage}

The common support header provides a growable byte buffer with checked growth, explicit little-endian field access, file I/O, and a bounded allocation registry.
All tool-owned allocations live until process exit and are released by one registered cleanup function.
The Lua state is also closed there.
This avoids mixing temporary symbol-string ownership, object-buffer ownership, and relocation ownership in a small command-line program.
It is not a recommended lifetime model for a long-running compiler server, which would need a per-request context and reusable error returns.

Buffers have a 16 MiB cap, symbols a 4096-entry cap, and relocations an 8192-entry cap.
The caps make resource failures explicit and bound the review examples.
They do not change the meaning of a successfully assembled source.
A capacity failure is reported as a tool failure rather than silently dropping a symbol or truncating an output section.

Figure~\ref{fig:ch56-storage} shows the principal structures.
The source buffer is mutable while lines are tokenized.
Symbol names are separately copied, so a later split at a colon or comma cannot change a previously interned name.
Relocations retain numeric symbol identities rather than pointers into a resizable symbol array.

\begin{figure}[htbp]
\centering
\begin{tikzpicture}[node distance=16mm and 18mm]
\node[basic block] (text) {generated source buffer};
\node[basic block,below left=of text] (secs) {four section records};
\node[basic block,below right=of text] (sym) {stable symbol array};
\node[basic block,below=of sym] (rel) {fixup array};
\draw[cfg edge] (text)--(secs); \draw[cfg edge] (text)--(sym);
\draw[cfg edge] (rel)--node[right,font=\small] {symbol identifier}(sym);
\draw[cfg edge] (rel)--node[below,font=\small,fill=white,inner sep=2pt] {section + field offset}(secs);
\end{tikzpicture}
\caption{Stable identifiers connect fixups to symbols, while section records own the byte and memory extents.}
\label{fig:ch56-storage}
\end{figure}

\section{Run Lua before parsing any generated statement}
\label{sec:ch56-emit}

Listing~\ref{lst:ch56-emit} is the actual host callback.
It uses the explicit Lua string length, checks for an embedded source NUL, copies the text, and appends one newline.
The copy matters because the pointer returned by Lua refers to Lua-managed storage.
The assembler's source stream must remain valid independently of later Lua garbage collection.

\begin{listing}[htbp]
\inputminted[firstline=14,lastline=20]{c}{example-code/lua-toolchain/luaas.c}
\caption{The emit callback copies Lua bytes into the assembler-owned source stream.}
\label{lst:ch56-emit}
\end{listing}

The script may emit several lines at once or one line at a time.
Their concatenation is parsed only after the protected Lua call returns successfully.
Consequently, an error after an earlier \code{emit} call still prevents assembly and object publication.
The test suite explicitly executes a script that emits a no-operation and then raises an error, verifying that no object is produced.
The Lua API and its protected-call behavior are specified by the reference manual~\cite{ierusalimschy-2026-lua}.

The callback does not assign target addresses.
It also does not evaluate \code{.quad answer} as a Lua variable access.
Those bytes describe a later assembler statement, preserving the separation between the generator's variables and the target program's symbols.

\section{Parse statements and bind labels}
\label{sec:ch56-parser}

For each generated line, the parser removes its semicolon comment, trims whitespace, recognizes an optional label, and separates opcode/directive from operands.
A label binds the current section offset before the rest of that line emits bytes.
Section directives only select the active counter; they do not reset it.
Binding directives change the symbol's linkage attribute without inventing a definition.

The parser is intentionally not a general token framework.
Its line language has no quoted string directive, so semicolon handling cannot accidentally remove a semicolon inside such a string.
Lua emits target string data as numeric bytes instead.
If a future version adds quoted assembly strings, comment scanning must become quote-aware before that syntax is accepted.
This is an example of keeping the parser's implementation assumptions aligned with its advertised grammar.

Expression parsing accepts a signed integer or a symbol with an optional integer addend.
Arithmetic overflow is checked before negating or adjusting the addend.
A branch or RIP-relative address-forming instruction requires a symbolic expression so that the place-relative adjustment is unambiguous.
Raw numeric displacement syntax could be added, but would need to state whether its number is a displacement or an absolute target.

\begin{algorithm}[htbp]
\caption{Assemble one line under the case-study grammar.}
\label{alg:ch56-line}
\begin{pseudocode}
\Require one generated line and current section/symbol state
\Ensure state update, exact bytes/fixup, or a line-specific error
\State remove comment and surrounding whitespace; ignore an empty line
\If{a label precedes the statement}
  \State validate its spelling and require no previous definition
  \State define it at the current section-relative offset
\EndIf
\State classify the remaining operation as directive or instruction
\If{directive}
  \State validate all operands and update section, binding, alignment, or storage
\Else
  \State validate operand classes and select a supported fixed-size form
  \State emit known bytes and retain symbolic fields as fixups
\EndIf
\end{pseudocode}
\end{algorithm}

Algorithm~\ref{alg:ch56-line} is the organizing control flow; the source contains the concrete checks.
The most useful tests target decisions that could otherwise be silently coerced: incompatible register widths, an unsupported register, an oversized immediate, a symbolic byte, and initialized bytes inside NOBITS storage.
Each must produce a diagnostic rather than an object that happens to have a valid ELF header.

\section{Reserve fields without losing relocation meaning}
\label{sec:ch56-fields}

The field-emission routine receives a classified expression, width, and relocation type.
For constants it verifies the field's permitted range and writes the low bytes in little-endian order.
For a symbol it creates a fixup at the current section offset, then writes exactly the required number of zero placeholder bytes.
For PC32 fields it subtracts four from the expression addend because the supported displacement fields end at the instruction boundary.
This adjustment is validated against signed-64 overflow before the subtraction.

No later code guesses where the field begins by decoding the instruction again.
The exact field offset is captured at reservation time.
Figure~\ref{fig:ch56-fixup} shows the call in the supplied demo object.
Its opcode is at text offset 17 and its fixup at offset 18; the later linker places the target at text-base plus 32.

\begin{figure}[htbp]
\centering
\begin{tikzpicture}[node distance=18mm]
\node[basic block] (bytes) {text+17: E8 00 00 00 00};
\node[basic block,below=of bytes] (fix) {offset 18; PC32; answer; addend -4};
\node[basic block,below=of fix] (sym) {answer: global, undefined};
\draw[cfg edge] (fix)--node[right,font=\small] {patch four bytes}(bytes);
\draw[cfg edge] (fix)--node[right,font=\small] {stable identity}(sym);
\end{tikzpicture}
\caption{The placeholder bytes and relocation record are emitted together, before the target address is known.}
\label{fig:ch56-fixup}
\end{figure}

\section{Walk the complete demo object}
\label{sec:ch56-demo}

The demo initializes \code{ecx} to three, decrements it in a loop, calls \code{answer}, moves the returned value to the exit-status register, and invokes the Linux exit system call.
The separately assembled \path{answer.lua} supplies a function returning $42$.
The first object also contains a pointer initialized to \code{answer} and a 32-byte zero-filled scratch area.
Together these exercise local PC-relative control flow, an external call, absolute pointer relocation, section placement, and NOBITS storage.

Table~\ref{tab:ch56-offsets} computes the first object's text layout.
The subtract uses the fixed six-byte \code{81 /5 imm32} form rather than a shorter immediate encoding.
The conditional branch uses the six-byte near form.
These choices explain the offsets; they are not estimates based on source line count.

\begin{table}[htbp]
\caption{The demo's exact fixed forms produce a 31-byte text contribution.}
\label{tab:ch56-offsets}
\centering
\begin{tabular}{llll}
\toprule
Operation & Start & Length & Next offset \\
\midrule
Initialize ecx & $0$ & $5$ & $5$ \\
Subtract one & $5$ & $6$ & $11$ \\
Conditional branch & $11$ & $6$ & $17$ \\
Call answer & $17$ & $5$ & $22$ \\
Move eax to edi & $22$ & $2$ & $24$ \\
Load exit syscall number & $24$ & $5$ & $29$ \\
System call & $29$ & $2$ & $31$ \\
\bottomrule
\end{tabular}
\end{table}

The loop branch field begins at $13$ and targets offset $5$, so its final displacement is $5-4-13=-12$.
The external call's value remains unresolved in the object.
The data pointer contributes eight file bytes and an absolute relocation at data offset zero.
Scratch contributes zero payload bytes and a NOBITS extent of $32$ with alignment $16$.
Table~\ref{tab:ch56-fixups} records the three relocation relationships.

\begin{table}[htbp]
\caption{The demo carries three independent deferred computations across the object boundary.}
\label{tab:ch56-fixups}
\centering
\begin{tabular}{llll}
\toprule
Target section & Field offset & Symbol & Type/addend \\
\midrule
Text & $13$ & private loop label & PC32, $-4$ \\
Text & $18$ & answer & PC32, $-4$ \\
Data & $0$ & answer & 64, $0$ \\
\bottomrule
\end{tabular}
\end{table}

\section{Serialize one coherent ELF graph}
\label{sec:ch56-object}

The writer creates four allocated sections, four corresponding RELA sections, a symbol table, a symbol string table, and a section-name string table, plus section zero.
Even an empty relocation section has correct type and cross-table links.
The symbol writer places local entries first and builds an internal-to-ELF index map before serializing fixups.
The file writer aligns each payload as required, records its actual offset, and writes the section-header table last.

\begin{algorithm}[htbp]
\caption{Finish the case-study relocatable object only after all consistency checks succeed.}
\label{alg:ch56-object}
\begin{pseudocode}
\Require section buffers, symbol entries, and stable-identifier fixups
\Ensure a self-consistent ELF64 ET\_REL file
\State reject unresolved names lacking permitted external binding
\State build string tables and local-first symbol serialization
\State translate every fixup's target to its final ELF symbol index
\State append aligned section payloads, excluding NOBITS contents
\State append relocation and symbol/string payloads
\State append section headers containing the recorded offsets and links
\State complete the ELF header with section-table location and sizes
\State publish the complete file
\end{pseudocode}
\end{algorithm}

Figure~\ref{fig:ch56-file} emphasizes that NOBITS size is metadata rather than an omitted chunk that still occupies file space.
Its section header records the logical extent.
The loader does not read that extent from the relocatable object; the later linker incorporates it into the memory size of a loadable segment.
Confusing these two file types is a common source of incorrect BSS handling.

\begin{figure}[htbp]
\centering
\begin{tikzpicture}[node distance=10mm]
\node[basic block] (h) {ELF64 header};
\node[basic block,below=of h] (p) {text / data / relocation payloads};
\node[basic block,below=of p] (s) {symbols and string tables};
\node[basic block,below=of s] (t) {section headers, including BSS size};
\draw[cfg edge] (h.east) to[out=0,in=0] node[right,font=\small] {section-table offset}(t.east);
\draw[cfg edge] (t.west) to[out=180,in=180] node[left,font=\small] {payload offsets}(p.west);
\end{tikzpicture}
\caption{The writer records BSS extent in metadata while keeping the corresponding initialized payload absent.}
\label{fig:ch56-file}
\end{figure}

\section{Build and verify through independent paths}
\label{sec:ch56-tests}

The build script requires a C11 compiler, Lua 5.4 headers/library, and \code{pkg-config}.
Its output directory is supplied explicitly so generated executables remain outside the source tree.
The following commands assemble the two demo inputs:
\begin{verbatim}
sh example-code/lua-toolchain/build.sh /tmp/opencode/lua-tools gcc
/tmp/opencode/lua-tools/luaas example-code/lua-toolchain/demo.lua \
  /tmp/opencode/lua-tools/demo.o
/tmp/opencode/lua-tools/luaas example-code/lua-toolchain/answer.lua \
  /tmp/opencode/lua-tools/answer.o
\end{verbatim}

The test driver compiles both tools with GCC and Clang, normally and with address and undefined-behavior sanitizers.
It compares selected encodings against GNU \code{as}, checks object metadata with \code{readelf}, links the objects with GNU \code{ld}, and executes the native program.
Chapter~\ref{ch:lua-linker} adds the independent layout and relocation checks for our linker.
Whole-object bytes are compared across configurations to detect nondeterministic ordering or uninitialized serialized padding.

Negative tests are equally important.
Duplicate labels, undeclared external references, unsupported registers, bad alignment, oversized data, mismatched operand widths, and Lua errors must all fail before output publication.
The tests also use an error object that is a Lua table, exercising the non-string error path.
These checks establish concrete behavior on their tested cases; they do not establish correctness for arbitrary Lua programs or every possible assembly file.

\section{Exercises}
\label{sec:ch56-exercises}
\begin{enumerate}
\item\label{ex:ch56-field} Starting from Table~\ref{tab:ch56-offsets}, derive both PC32 field offsets without reading the relocation table.
\item\label{ex:ch56-counter} Change the countdown's runtime count from three to seven. Which instruction bytes change, and which symbol offsets remain unchanged?
\item\label{ex:ch56-size} Compare the demo's eight-byte initialized pointer and 32-byte scratch area. State their object payload sizes, memory extents, and relocation needs.
\item\label{ex:ch56-script-error} Explain why executing the script and assembling each emitted line immediately would complicate the guarantee that a later Lua error produces no object.
\item\label{ex:ch56-differential} Why does the reference encoding test request \code{movabs} for a full-width 64-bit immediate move instead of accepting the production assembler's default choice?
\end{enumerate}

\section{Further reading}
\label{sec:ch56-reading}
The Lua reference manual defines the embedding operations used by the host~\cite{ierusalimschy-2026-lua}.
The generic ELF and AMD64 specifications define the serialized contract independently of our C implementation~\cite{xinuos-2026-elf,matz-2012-amd64}.
Salomon provides broader context for assembler listings, diagnostics, and implementation choices~\cite{salomon-1993-assemblers}.
